%\documentclass[11pt,letterpaper]{article}
\documentclass[[a4paper,landscape]{article}\usepackage[]{graphicx}\usepackage[]{color}
% maxwidth is the original width if it is less than linewidth
% otherwise use linewidth (to make sure the graphics do not exceed the margin)
\makeatletter
\def\maxwidth{ %
  \ifdim\Gin@nat@width>\linewidth
    \linewidth
  \else
    \Gin@nat@width
  \fi
}
\makeatother

\definecolor{fgcolor}{rgb}{0.345, 0.345, 0.345}
\newcommand{\hlnum}[1]{\textcolor[rgb]{0.686,0.059,0.569}{#1}}%
\newcommand{\hlstr}[1]{\textcolor[rgb]{0.192,0.494,0.8}{#1}}%
\newcommand{\hlcom}[1]{\textcolor[rgb]{0.678,0.584,0.686}{\textit{#1}}}%
\newcommand{\hlopt}[1]{\textcolor[rgb]{0,0,0}{#1}}%
\newcommand{\hlstd}[1]{\textcolor[rgb]{0.345,0.345,0.345}{#1}}%
\newcommand{\hlkwa}[1]{\textcolor[rgb]{0.161,0.373,0.58}{\textbf{#1}}}%
\newcommand{\hlkwb}[1]{\textcolor[rgb]{0.69,0.353,0.396}{#1}}%
\newcommand{\hlkwc}[1]{\textcolor[rgb]{0.333,0.667,0.333}{#1}}%
\newcommand{\hlkwd}[1]{\textcolor[rgb]{0.737,0.353,0.396}{\textbf{#1}}}%
\let\hlipl\hlkwb

\usepackage{framed}
\makeatletter
\newenvironment{kframe}{%
 \def\at@end@of@kframe{}%
 \ifinner\ifhmode%
  \def\at@end@of@kframe{\end{minipage}}%
  \begin{minipage}{\columnwidth}%
 \fi\fi%
 \def\FrameCommand##1{\hskip\@totalleftmargin \hskip-\fboxsep
 \colorbox{shadecolor}{##1}\hskip-\fboxsep
     % There is no \\@totalrightmargin, so:
     \hskip-\linewidth \hskip-\@totalleftmargin \hskip\columnwidth}%
 \MakeFramed {\advance\hsize-\width
   \@totalleftmargin\z@ \linewidth\hsize
   \@setminipage}}%
 {\par\unskip\endMakeFramed%
 \at@end@of@kframe}
\makeatother

\definecolor{shadecolor}{rgb}{.97, .97, .97}
\definecolor{messagecolor}{rgb}{0, 0, 0}
\definecolor{warningcolor}{rgb}{1, 0, 1}
\definecolor{errorcolor}{rgb}{1, 0, 0}
\newenvironment{knitrout}{}{} % an empty environment to be redefined in TeX

\usepackage{alltt}
\usepackage{datetime}
\usepackage[utf8]{inputenc}
\usepackage[spanish]{babel}
\usepackage{graphicx}
\usepackage{fancybox}
\usepackage{fancyhdr} %pie de pagina y cabecera
\usepackage[utf8]{inputenc} %solucion del problema de los acentos.
\usepackage{multicol,pst-plot}
\usepackage[left=1.5cm,right=1.5cm,top=2cm,bottom=2cm]{geometry}
\usepackage{xcolor}
\usepackage[many]{tcolorbox}
\usepackage{svg}
\usepackage{float}
\usepackage{wrapfig}

%%%%%%%%%  ESTILO DOC %%%%%%%%%%%%%%%%%%%%%%%%%%%%%%%%%%%%
\definecolor{meteoblue}{RGB}{0,160,250}

\tcbset{mystyle/.style={
  breakable,
  enhanced,
  outer arc=0pt,
  arc=0pt,
  colframe=meteoblue,
  colback=meteoblue!20,
  attach boxed title to top left,
  boxed title style={
    colback=meteoblue,
    outer arc=0pt,
    arc=0pt,
    },
  title=Example~\thetcbcounter,
  fonttitle=\sffamily
  }
}

\newtcolorbox[auto counter,number within=section]{example}[1][]{
  mystyle,
  title=NOTA,
  }





%%%%%%%%%%%%%%%%%%%%%%%%%%%%%%%%%%%%%%%%%%%%%%%%%%
\pagestyle{fancy}
\fancyhead[L]{
\textcolor{meteoblue}{
\bfseries Alpred S.L. \\
\small Servicio Metorológico a Empresas
}
}
\fancyhead[R]{\includegraphics[width=2.5cm]{/home/raquel/repos/pruebas/informes/meteoredlogo.pdf}}
\fancyhead[C]{ Previsi\'on para {Quito [-0.2300;281.4800]} \\
\today}
%\fancyfoot{}
\fancyfoot[R]{Tel. 968485304 Fax. 968440315}
\fancyfoot[L]{C/ Mayor, 46 \\
30893 - Almendricos (Murcia)
}
\fancyfoot[C]{\today}
%%%%%%%%%%%%%%%%%%%%%%%%%%%%%%%%%%%%%%%%%%%%%%%%%%
\IfFileExists{upquote.sty}{\usepackage{upquote}}{}
\begin{document}
%\begin{center}\vspace{-1cm}
%\textcolor{meteoblue}{\textbf{ \huge Predicci\'on para \textcolor{meteoblue}{result$nombre} }}\\
%\date{\today}
%\end{center}
%\textcolor{meteoblue}{\rule{\linewidth}{1mm}}

\begin{figure}
\begin{knitrout}
\definecolor{shadecolor}{rgb}{0.969, 0.969, 0.969}\color{fgcolor}
\includegraphics[width=\maxwidth]{figure/Figtemperatura-1} 

\end{knitrout}
\end{figure}


\begin{figure}
\begin{knitrout}
\definecolor{shadecolor}{rgb}{0.969, 0.969, 0.969}\color{fgcolor}
\includegraphics[width=\maxwidth]{figure/Figprecipitacion-1} 

\end{knitrout}
\end{figure}


\begin{figure}
\begin{knitrout}
\definecolor{shadecolor}{rgb}{0.969, 0.969, 0.969}\color{fgcolor}
\includegraphics[width=\maxwidth]{figure/Figdir2-1} 

\end{knitrout}
\end{figure}

\clearpage
\newpage

\textbf{\textcolor{meteoblue}{INFORMACI\'ON}}\\

Descripción de las variables meteorol\'ogicas.

\begin{itemize}
    \item \textbf{Temperatura m\'axima y m\'inima}: Es la temperatura máxima y mínima diaria que tiene luegar entre las 00 h y las 23:59 h.
     Viene dada en grados centigrados ($^\circ$C)

     \item \textbf{Punto de roc\'io}: Temperatura a la que el aire, a presi\'on y contenido de humedad constantes,
     se satura (se condensa).  Su unidad de medida es la misma que la temperatura, es decir,
     grados centigrados ($^\circ$C).

     \item \textbf{Humedad relativa}: Se expresa en tanto por ciento.
     Cuanto mayor es el porcentaje más cerca se está de la saturación (formación de pequeñas gotas de agua).

    \item \textbf{Presi\'on atmosf\'erica}: Es la presión reducida a nivel del mar.
    La presión disminuye exponencialmente con la altura (aproximadamente 1hPa por cada 8 m de ascenso).
    La unidad del sistema internacional es el Pascal. También se expresa en milibares (1 mbar = 1 hPa). La presi\'on normal es de 1013,2 mbar.
    Las altas presiones se asocian a anticiclones y las bajas presiones a borrascas. No siempre el mal tiempo va asociado a bajas temperaturas.

    \item \textbf{Precipitaci\'on}: Cantidad de agua que precipita desde las nubes y llega al suelo en distintas formas: lluvia, nieve,  granizo
    Viene definida por la suma total de la precipitaci\'on estratiforme
    (tipo de precipitaci\'on asociada a la nubosidad a gran escala, generalmente asociada a sistemas frontales)
    y la convectiva (precipitaci\'on relacionada con la nubosidad de crecimiento vertical, que en gran porcentaje da lugar a chubascos y tormentas).
    La precipitaci\'on se mide en mil\'imetros de espesor acumulado en altura (mm),
    que es equivalente a la medida de litros por metro cuadrado (l/m$^2$).\\

    \item \textbf{Nubosidad}: Se mide en porcentaje (semiesfera).
    Un valor de 0\% significa un cielo totalmente despejado duarnate un periodo de tiempo determinado.
    Un cielo totalmente cubierto conllevaría un valor de 100\%.

    \item \textbf{Radiaci\'on Solar}: Cantidad de energía solar que incide por unidad de área horizontal durante un periodo determinado (en este informe cada 3 h)
    y se mide en Julios/m$^2$. NOTA: no se si es directa + indirecta.

    \item \textbf{Direcci\'on del viento}: La dirección del viento nos indica de donde viene éste.
     La dirección viene marcada por la rosa de los vientos; situando el norte a 0$^\circ$
    y recorriendo en sentido horario se sitúa el este a  90$^\circ$, el sur a 180 $^\circ$ y el oeste a 270 $^\circ$.

    \item \textbf{Velocidad del viento}: Magnitud cuya unidad mas común es el kilómetro por hora (km/h).

\end{itemize}


%NOTA: La fiabilidad de las predicciones disminuye a lo largo del tiempo. METEORED ofrece sus predicciones permanente en su web (www.tiempo.com)
%IFS (Integrated Forecasting System)

\vspace{1cm}
\begin{example}
La magnitud de cada variable se proporciona cada tres horas. Todas las variables tienen un valor cada tres horas,
excepto la precipitaci\'on y la radiaci\'on solar que son acumulables en periodos de tres horas.\\
%La temperaturea maxima es max(temperatura) \\
%La temperaturea media es mean(temperatura)
\end{example}

\clearpage
\newpage




\begin{center}
    \includegraphics[width=\textwidth]{./figure/Figresumen-1.pdf}
\end{center}


\begin{footnotesize}
\textbf{\textcolor{meteoblue}{INFORMAC\'ION}}
\begin{itemize}
    \item \textbf{Presi\'on atmosf\'erica}: Está referida al nivel del mar (mbar).
    \item \textbf{Precipitaci\'on}: estratiforme (tipo de precipitaci\'on asociada a sistemas frontales)
            y la convectiva (nubosidad de crecimiento vertical), (mm).
    \item \textbf{Nubosidad}: Se mide en porcentaje (semiesfera).
    \item \textbf{Temperatura m\'axima y m\'inima}: Viene dada en grados centigrados ($^\circ$C)
    \item \textbf{Punto de roc\'io}: Temperatura a la que el aire, a presi\'on y contenido de humedad constantes,
            se satura (se condensa). ($^\circ$C).
    \item \textbf{Direcci\'on del viento}: Esquema de direcci\'on del viento.
    \item \textbf{Velocidad del viento}:(km/h).
    \item \textbf{Radiaci\'on Solar}: Cantidad de energía acumulada en un periodo de tiempo y se mide en megajulios/m$^2$
\end{itemize}
\end{footnotesize}



%\begin{figure}
%<<Figdir3, echo=FALSE, fig.width=10, results = 'hide'>>=
%dias=date(kk$dateok)
%horas=hour(kk$dateok)
%kk <- within(kk, {
%dirwind <- -dir +90
%x.end <- dias + 0.5*cos (dirwind/180 * pi)
%y.end <- as.numeric(horas) + 0.5*sin (dirwind/180 * pi)
%dia <- dias
%hora <- horas
%})

%ggdir3 <- ggplot(kk, aes(x= dia, y=hora , fill=wind))+
%            geom_raster() +
%            scale_fill_distiller(palette="Blues",direction=1,breaks=seq(0,20,2),labels=seq(0,20,2),guide = "legend",name="Viento (m/s)")+
%            geom_segment(aes(xend= dia, x= x.end, yend = hora, y = y.end), arrow = arrow(length = unit(0.2,"cm")))+
 %          # geom_text(aes(label=rumbo, size=10), color="gold")+
 %           scale_x_date(breaks = date_breaks("1 days"),date_labels = "%b \n %a\n %d")+
 %           scale_y_continuous(breaks=seq(2,23,3))+
 %         #  scale_y_continuous(trans = "reverse",breaks=seq(2,23,3))+
 %           theme(legend.position = 'bottom', legend.direction = 'horizontal', legend.text = element_text(colour="black", size=10))+
 %           guides(fill= guide_legend(nrow=1,barheight=0.4))+
 %           labs(title="Dirección y velocidad del viento a 10 m",y="Hora",x="")
%#axis.title.x=element_blank(), axis.text.x=element_blank(), axis.ticks.x=element_blank()

%ggdir3
%@
%\end{figure}


%<<CreateTable1, warnings=FALSE ,echo=FALSE, results = 'asis'>>=
%kktabla=cbind(kk$dateok,kk$presion,kk$nub,kk$hum,kk$precip)
%table_latex <- xtable(kktabla,caption = "cosas")
%print(table_latex,
%   latex.environments =  c("scriptsize", "center", "widestuff"),
%	          floating = FALSE
%		  )
%@


%\begin{figure}
%<<Figdir, echo=FALSE>>=
%tstep=seq(1,16)
%par(mfrow=c(4,4))
%for (i in tstep) {
%    plot(c(0,kk6h$u[i]),c(0,kk6h$v[i]),xlim=c(-1,1),ylim=c(-1,1),type="l", ylab="", xlab="",xaxt= "n", yaxt="n", main=paste(kk6h$dir[i]," H=", kk6h$date[i] ,sep=""),col="red")
%    axis(2, at=c(-1,0,1),labels=c("SW","W","NW"), col.axis="red", las=2)
%    axis(4, at=c(-1,0,1),labels=c("SE","E","NE"), col.axis="red", las=2)
%    axis(1, at=c(0),labels=c("S"), col.axis="red", las=2)
%    axis(3, at=c(0),labels=c("N"), col.axis="red", las=2)
%}
%@
%\end{figure}

\vspace{5cm}
\newpage

%<<CreateTable,  warnings=FALSE, echo=FALSE, results = 'asis'>>=
%kk=cbind(fecha[t3h], presion[t3h],temperatura[t3h], dir[t3h])
%table_latex <- xtable(kk6h,caption = "cosas")
%print(table_latex,
%   latex.environments =  c("scriptsize", "center", "widestuff"),
%	          floating = FALSE
%		  )
%@




\end{document}
